% Requires: \usepackage{amsmath,amssymb,booktabs,tabularx}
\providecommand{\Ftgt}{F^{\star}}
\providecommand{\Fmin}{F_{\min}}
\providecommand{\Fmax}{F_{\max}}
\providecommand{\Faud}{F^{\mathrm{aud}}}
\providecommand{\ind}{\mathbb{I}}

\section{Simulation Task and Evaluation Suite}
\label{sec:benchmark}

We use ForceWipe as a controlled simulation testbed for direct first-pass
wiping. A successful evaluation must traverse the path, maintain useful
contact, track the target normal force, and avoid recorded force-limit
exceedances. The test cases vary path geometry, surface orientation, tool
profile, contact parameters, disturbances, and weak surface curvature. The
resulting suites define controlled and transfer conditions for the same direct
action interface.

\subsection{Direct First-Pass Task}
\label{subsec:task}

A seven-degree-of-freedom Franka Emika Panda manipulates a wiping tool in
ManiSkill3 with the PhysX central processing unit (CPU)
backend~\cite{xiang2020sapien,c_maniskill}.
Physics and control advance together at $100$~Hz. At each step, the local task
frame is formed by the path tangent $\mathbf{t}_t$, outward surface normal
$\mathbf{n}_t$, and cross-track axis
$\mathbf{c}_t=\mathbf{n}_t\times\mathbf{t}_t$. The controller outputs
\begin{equation}
\mathbf{a}_t=
\begin{bmatrix}a_t^{\parallel}&a_t^{\perp}&a_t^{n}\end{bmatrix}^{\!\top}
\in[-1,1]^3,
\label{eq:direct_action}
\end{equation}
where the components request along-path, cross-track, and inward-normal
displacements. The fixed, force-independent actuator map is
\begin{equation}
\Delta\mathbf{x}_t=
(0.50\,\mathrm{mm})a_t^{\parallel}\mathbf{t}_t
+(0.25\,\mathrm{mm})a_t^{\perp}\mathbf{c}_t
-(0.50\,\mathrm{mm})a_t^{n}\mathbf{n}_t.
\label{eq:task_frame_map}
\end{equation}
Post-policy processing consists exclusively of action-box clipping and a
force-independent coordinate transform in \eqref{eq:task_frame_map}.

The tangent and normal are evaluated from the registered nominal path and
surface model at the simulator-state tool pose. They are not inferred from a
camera, force sensor, or the TD-MPC2 latent state. Consequently, ``direct'' in
this study denotes the absence of a separate force-control or shielding layer;
it does not denote model-free geometric perception. This registered-frame
assumption is shared by all compared methods.

At decision step $t$, $F_t$ denotes the force supplied to the causal
observation before $\mathbf{a}_t$ is selected. We reserve $\Faud_t$ for the
native force read back after the corresponding physics step and attached to
that transition for evaluation. The target normal force is
$\Ftgt\in\{5,8,12\}$~N. The online contact indicator uses
$F_t\geq\Fmin=3$~N, whereas evaluation metrics use the post-step contact set
$\Faud_t\geq\Fmin$ defined in Sec.~\ref{subsec:metrics}. A strict sampled violation is
$\Faud_t>\Fmax=15$~N. The statement ``zero
violations'' therefore means that no recorded $100$-Hz sample exceeded
$15$~N. All force-limit results use this native sampled definition.
The three force targets, the 3-N contact threshold, and the 15-N sampled
limit form the frozen operating contract of this simulation task. They are
used identically by every method and were fixed before the controlled and
factor-separated evaluations; they are not presented as general robot or
human-contact safety standards.

Each path is divided into 20 normalized progress bins. A bin is marked
complete once it has accumulated at least two samples satisfying
$\Fmin\leq\Faud_t\leq\Fmax$, corresponding to $20$~ms at the simulator
rate. A task pass requires a single first-pass lifecycle to terminate with
all 20 bins complete and progress $s\geq0.98$. Re-wiping is disabled in all
three evaluation suites. The dose endpoint is synthetic contact
coverage and serves as the completion proxy.

\subsection{Controlled Evaluation Blocks}
\label{subsec:controlled_suite}

The initial and held-out controlled suites use three
surface orientations (flat, $5^{\circ}$ incline,
and $10^{\circ}$ incline), five path families (line, diagonal, arc,
S-curve, and polyline), and three tool profiles (wide-soft, narrow-soft, and
wide-hard). Their synthetic residual fields comprise single-blob, islands,
edge-dense, stripe or multimodal, and Gaussian-random-field patterns. These
factors are assembled into prespecified blocks rather than sampled as an
independent factorial design.

The five initial blocks Q0--Q4 use the nominal contact plant and
no exogenous disturbance. They vary path length between approximately
$0.156$ and $0.166$~m while changing geometry, tool, and residual field by
block. The ten previously unused held-out blocks C0--C9 form two disjoint
groups. C0--C4 repeat the five path families on a nominal plant. C5--C9
pair the same roster positions with randomised path length, friction,
support stiffness, effective mass, damping, and restitution. Their frozen
definitions also label designated blocks with lateral or normal impulses,
reference noise, and delayed sensor noise. A later binding audit established
the physical impulse call chain but did not establish that the declared
reference-noise and sensor-noise signals reached the direct controller's
learner-visible force path. We therefore make no robustness claim for those
two declared noise conditions. Because several verified attributes change
jointly, grouped path, surface, tool, residual, and plant summaries quantify
block-level transfer rather than single-factor effects.

\subsection{Weak-Curvature Out-of-Distribution Blocks}
\label{subsec:weak_curvature_suite}

The matched OOD study contains six previously unused, prespecified blocks
F0--F5. All surfaces are static cylinders with radii $1.2$--$1.8$~m. Its paths
are arc, S-curve, or polyline trajectories with lengths $0.13$--$0.205$~m.
Across the six blocks, the study includes all three tool profiles, six
residual families, randomised friction and contact plant parameters, and
predefined disturbance labels. Physical lateral and normal impulse call chains
are present in their designated blocks; the declared reference-noise and
sensor-noise/latency effects are excluded from the evidence claims because
their runtime binding was not established. Table~\ref{tab:evaluation_suites}
summarizes the evaluation stages.

\begin{table*}[t]
\centering
\caption{Prespecified simulation suites. Evaluation totals are Cartesian
products of blocks, force targets, and method replicates; blocks are the
environment units. For the deployment-gap panel, 375+45 denotes new stress
evaluations plus reused nominal evaluations.}
\label{tab:evaluation_suites}
\small
\setlength{\tabcolsep}{4.5pt}
\begin{tabular}{@{}lccccp{6.1cm}@{}}
\toprule
Stage & Blocks & Targets & Replicates & Evaluations & Main variation \\
\midrule
Controlled evaluation
& 5 & 3 & 5 checkpoints & 75
& Five paths; flat/$5^{\circ}$/$10^{\circ}$ surfaces; three tools; nominal plant \\
Held-out controlled evaluation
& 10 & 3 & 5 checkpoints & 150
& Five nominal multipath blocks plus five randomised-plant/disturbance blocks \\
Matched weak-curvature OOD
& 6 & 3 & method dependent & 288
& Static weak-curvature cylinders, three paths, three tools, randomised plant, and disturbances \\
Numerical sentinel panel
& 4 & 1 & one checkpoint & 36
& Four 12~N OOD baseline reproductions plus 32 solver, contact-offset, mesh, and interaction cells \\
Deployment-gap stress
& 3 & 3 & 5 checkpoints & 375+45
& Learner-visible force noise/bias, delay, TCP-position error, and registered-frame normal error \\
\bottomrule
\end{tabular}
\end{table*}

The controlled and weak-curvature suites are reported separately to preserve
their deliberately different scenario distributions.

\subsection{Evaluation Metrics and Success Criteria}
\label{subsec:metrics}

Force tracking is computed over the finite contact set
\begin{equation}
\mathcal{T}_{c}=\{t\mid \Faud_t\geq3~\mathrm{N}\}.
\end{equation}
If $\mathcal{T}_{c}$ is empty, the force-tracking metrics are undefined and
the evaluation fails the tracking criterion. For a nonempty contact set, we
report mean relative error (MRE), root-mean-square error (RMSE), and normalized
root-mean-square error (NRMSE):
\begin{align}
\overline F_c
&=\frac{1}{|\mathcal{T}_{c}|}
\sum_{t\in\mathcal{T}_{c}}\Faud_t,\nonumber\\
b_F&=\overline F_c-\Ftgt,\qquad
\mathrm{MRE}=\frac{|b_F|}{\Ftgt},\nonumber\\
\mathrm{RMSE}
&=\sqrt{\frac{1}{|\mathcal{T}_{c}|}
\sum_{t\in\mathcal{T}_{c}}(\Faud_t-\Ftgt)^2},\nonumber\\
\mathrm{NRMSE}&=\frac{\mathrm{RMSE}}{\Ftgt}.
\label{eq:force_metrics}
\end{align}
Tracking is considered successful when
$\mathrm{MRE}\leq0.15$ and $\mathrm{NRMSE}\leq0.20$. Signed bias is
retained so systematic under-tracking remains visible even when both
tolerances are met. These two tolerances are likewise study-defined acceptance
criteria frozen with the evaluation protocol, rather than thresholds inferred
from the reported trials.

The force evaluation records episode peak force, peak overshoot
$\max(\max_t\Faud_t-\Ftgt,0)$, peak-force mean, 95th percentile (P95), and
maximum, and sample counts
above $14.5$ and $15$~N. The $14.5$~N count is descriptive headroom; the
sampled-force-limit criterion requires zero samples above $15$~N. Contact metrics
include contact fraction and, after first contact, the number, total duration,
and maximum duration of maximal consecutive intervals below $3$~N.
Force-tail quantities describe the frozen numerical configuration in which
they were observed. Section~\ref{sec:credibility_audit} separately tests their
sensitivity to selected simulator settings.

An evaluation satisfies deployment integrity when every trace confirms
direct-action interface compliance: the selected learned branch issued the
executed action, followed by the fixed map in \eqref{eq:task_frame_map}. A
\emph{compound pass} is the intersection of task completion, tracking,
sampled-force-limit compliance, and deployment integrity. Nonfinite force or
an empty contact set counts as an evaluation failure. Planning time and the
actor/MPPI action fractions are descriptive computational metrics.

\subsection{Scope of the Suite}
\label{subsec:scope}

The blocks are fixed conditions, with each block defining an
environment unit. Replicated checkpoint--block--target trials increase the
precision of within-design comparisons. The suite characterizes direct
force-aware wiping under controlled variation and weak-curvature transfer;
Sec.~\ref{sec:discussion} defines the corresponding deployment scope.
